%
%
\section*{Framework introduction}\addcontentsline{toc}{section}{Framework introduction}

A positive integer is \emph{powerful} if every prime that divides it divides it at least twice. Two consecutive integers can both be
powerful, as $8, 9$ shows, and there are infinitely many such pairs. No three consecutive powerful numbers are known. We expect that there
are none; the introduction of Part~I gives the reasons, and none of them is a proof. This volume collects three parts on these two
situations.

\emph{Part~I} concerns triples $n - 1, n, n + 1$. It translates the middle $n$ into an address $(m, k)$ on a Pell equation, bounds $n$ from
below in terms of the kernel $m$, and profiles a counterexample: where it would have to live, and which kernels are the weak point of the
bound. \emph{Part~II} concerns pairs $n, n + 1$. It enumerates the pairs by families, gives an explicit lower bound for their count, compares
the count with a heuristic yardstick, and localizes the count of pairs at distance two in the kernels that Part~I singles out, together
with their counterparts of the other residue class.
\emph{Part~III} collects further observations: the gaps between consecutive powerful numbers, powerful values of cyclotomic polynomials,
and two tests of the independence assumptions behind the heuristic of Part~I, each with the power of the test.

\emph{How to read the volume.} The parts follow one arc. Part~I states the question and its address system, shows what can be excluded
and why the rest depends on witnesses, and names the step at which the methods stop; a certified witness is found for \Z{rgzert} of the
\Z{rgpaare} building blocks of its finite record. Part~II turns the two halves of a triple into pairs and asks how many there are: towers,
families, an explicit lower bound, and why local models mislead. Part~III collects what we saw around the problem: the gap spectrum,
cyclotomic values, measured independence, and negative results with their reach. The open problems at the end of this introduction are
what remains.

The parts share one source and one numbering: a statement numbered I.4.5 carries that number here and in Part~I. Every numbered
statement carries a label that says what kind of evidence stands behind it (proved, literature, computed, measured, heuristic, open); the
table at the end of each part lists them. A statement that combines several kinds carries the label of the part that needs the most trust: a proof that uses our own computation is labeled computed, and a proof that uses a published result is labeled proved. Every number that we computed is generated from the deposited data files or from a cited
constant; the build refuses a
number that has no source.

The author is an independent researcher who likes numbers and science, not a number theorist or mathematician by training. The series
therefore keeps every claim as narrow as its proof or computation allows and says plainly where, to our knowledge, a question remains
open. It grew out of about two months of research. Specialists are welcome to extend, correct, or cite this work; the author intends it to
help this field progress and hopes that a number theorist will take the next steps.
The question whether three consecutive powerful numbers exist can be settled in only two ways: a triple is found, or a general statement is found that holds at all of the infinitely many addresses at once. Such a statement would have to agree with all results in this field, the formulas and laws of this volume and of the literature, and those still to be found. It need not be a single law: it could also be one exclusion that covers all addresses, or many exclusions that together cover them, in the way that the exclusions for cube middles \cite{Chan2025,She2025,Sayim2026EMW} each cover a part. We do not know yet whether such a statement exists.

\section*{Results at a glance}\addcontentsline{toc}{section}{Results at a glance}

The table lists every statement of our own in each part, grouped by the type of evidence behind it, with the main statements in bold.
Statements taken from the literature are not listed here; the complete list of the statements of a part, with the type of each, stands at the
end of that part. The open questions are listed completely in the section on the open problems below.

\begingroup\small
\begin{longtable}{l l p{8.0cm}}
\hline
part & type & statements \\
\hline
\endhead
Part I (triples) & proved & \textbf{I.2.1~height inequality}\newline \textbf{I.2.4~two wings}\newline \textbf{I.4.6~reduction to a witness statement}\newline \textbf{I.4.8~the divisor count of the Wieferich conditions}\newline \textbf{I.6.10~Wieferich primes of the unit and the primitive parts}\newline \emph{also:} I.2.2~the constant in the height bound; I.2.3~kernel bound; I.3.2~weaker bounds with less work; I.4.3~valuation identity; I.4.4~exponent-3 condition; I.4.7~the $T_1$-part of a lattice point; I.4.10~what the witness condition does not follow from; I.4.11~where a witness in an outer number can sit; I.4.12~two elementary constraints, and the size of the gap; I.4.13~the cyclotomic axis \\
 & computed & \textbf{I.3.1~reduction}\newline \textbf{I.4.5~the $b$-box}\newline \textbf{I.5.1~the strip}\newline \textbf{I.5.3~smooth kernels}\newline \textbf{I.5.4~square-type middles}\newline \textbf{I.6.2~no triple below $\Z{ohnetripelsicher}$, and the pairs at distance two with middle divisible by four}\newline \textbf{I.6.8~the primitive part over an explicit finite set}\newline \textbf{I.6.13~a second primitive prime factor}\newline \textbf{I.7.1~profile}\newline \emph{also:} I.5.5~the share of the candidates covered; I.6.3~ladders for the known kernels dividing $U_1$; I.6.4~the smallest witnesses; I.6.5~the condition modulo $36$ and the witnesses; I.6.6~how many conditions there are, and how often they are met; I.6.7~middles that are powers of two; I.6.11~the Wieferich events of the units; I.6.12~Lucas--Wieferich primes; I.8.2~Yokoi's representation; I.8.3~small units \\
 & measured & \textbf{I.6.1~exponents at the rank} \\
 & open & \emph{finiteness of the kernels of wing (ii)} (Problem~\problemvon{I:q:walsh}); \emph{the primitive part of $T_d(m)$} (Problem~\problemvon{I:q:primindex}); \emph{a middle that is a power of $2$} (Problem~\problemvon{I:q:zweierpotenz}); \emph{Reinhart's condition (C)} (Problem~\problemvon{I:q:reinhart}) \\
\hline
Part II (pairs) & proved & \textbf{II.6.2~the count is a sum over towers}\newline \textbf{II.6.4~localization of the pair count}\newline \textbf{II.6.6~where Tao's argument is weakest}\newline \emph{also:} II.2.1~powerful numbers; II.6.5~what the localization says, and what it does not say; II.7.1~pairs on one Pell equation \\
 & computed & \textbf{II.3.1~the known pairs}\newline \textbf{II.4.1~explicit lower bound}\newline \textbf{II.5.2~the families with a square overtake the comparison function}\newline \textbf{II.6.1~the middles that are not divisible by four}\newline \textbf{II.6.3~lower bound at every height}\newline \textbf{II.7.2~the third number}\newline \emph{also:} II.3.2~completeness inside the families \\
 & heuristic & \textbf{II.5.1~comparison function} \\
 & open & \emph{convergence of the rates} (Problem~\problemvon{II:q:konvergenz}); \emph{growth of the tower rate} (Problem~\problemvon{II:q:wachstum}); \emph{the kernels that divide $U_1$} (Problem~\problemvon{II:q:ezwei}) \\
\hline
Part III (further observations) & proved & \textbf{III.2.2~two explicit chains of powerful values of $\Kreis{3}$ and $\Kreis{6}$}\newline \textbf{III.2.5~the simplest case of $\Kreis{8}$} \\
 & computed & \textbf{III.1.2~the head of the gap spectrum}\newline \textbf{III.2.3~powerful values up to a bound}\newline \emph{also:} III.1.3~a ceiling for common divisors; III.2.1~the values $x^2 + 1$; III.2.4~a class without squares \\
 & measured & \textbf{III.1.4~a counting model, measured}\newline \textbf{III.3.1~no joint event along the divisors, and the power of the test}\newline \textbf{III.4.1~no deviation from equidistribution, memory or coupling}\newline \emph{also:} III.1.5~common factors of pairs \\
 & heuristic & III.5.3~a conjecture on the most frequent gap \\
 & open & \emph{the most frequent gap} (Problem~\problemvon{III:q:luecken}); \emph{powerful values of $x^4 + 1$} (Problem~\problemvon{III:q:phiacht}) \\
\hline
\end{longtable}
\endgroup


\inwords A proved statement stands on a proof in the text, which may use cited results. A computed statement stands on a computation whose data are in the deposit; where it depends on a published search, the statement says so.
A measured statement describes a finite range and claims nothing beyond it, and a heuristic is a yardstick, not a claim.

\section*{Where this work meets known open problems}\addcontentsline{toc}{section}{Where this work meets known open problems}

The questions of this volume touch several known open problems. The table names each one, says what this volume reaches toward it,
and says what stays open. None of these problems is solved here.

\begingroup\small\renewcommand{\arraystretch}{1.25}
\begin{longtable}{>{\raggedright\arraybackslash}p{3.4cm} >{\raggedright\arraybackslash}p{6.9cm} >{\raggedright\arraybackslash}p{3.0cm}}
\hline
problem & what this volume reaches & what stays open \\
\hline
\endhead
Three consecutive powerful numbers (Erd\H{o}s; in a form about Pell numbers, Mollin and Walsh \cite[p.~111]{MollinWalsh1986a}) &
no triple has a middle below $\Z{ohnetripelsicher}$ (Proposition~\ref{I:prop:paarliste}(d); Table~\ref{tab:schranken} compares it with the
bounds in the literature, among them a machine-checked proof below $10^{14}$ \cite{MianSiddique2026}); every lattice point with kernel up to
$\Z{kernschranke}$ and middle below $10^{\Z{hoehe}}$ is excluded (Theorem~\ref{I:thm:reduktion}); the profile of a counterexample
(Theorem~\ref{I:thm:profil}) &
the question itself; Problems~\problemvon{I:q:walsh} to~\problemvon{I:q:reinhart} \\
Wieferich primes of Pell units (for primes that have a rank, the Lucas--Wieferich primes of McIntosh and Roettger \cite{McIntoshRoettger2007}) &
if the unit of a kernel has only finitely many Wieferich primes of odd rank, then all but finitely many of its building blocks are not
powerful (Proposition~\ref{I:prop:lucaswieferich}); below $\Z{wfp}$ the \Z{wfkerne} units of Proposition~\ref{I:prop:raenge} carry
\Z{wfereignisse} Wieferich events, \Z{wfungerade} of them of odd rank $d \ge 5$ (Remark~\ref{I:rem:wfdaten}) &
whether there are only finitely many; Problem~\problemvon{I:q:primindex} \\
Kernels $m$ with $m \mid U_1(m)$ (Walsh's question \cite[p.~97]{Walsh1988}) &
they are the weak point of the height bound (Corollary~\ref{I:cor:fluegel}); Reinhart's list of them is complete up to $1.5 \cdot 10^{12}$
\cite[Remark~5.5]{Reinhart2023}, and for the two known kernels $m \equiv 7 \pmod 8$ with $T_1$ even a ladder shows how large a middle would have to be
(Proposition~\ref{I:prop:reinhart}); the count of pairs at distance two is localized in them (Proposition~\ref{II:prop:lokalisierung}) &
finiteness, Problem~\problemvon{I:q:walsh}; a bound for $E_3$, Problem~\problemvon{II:q:ezwei} \\
The $abc$ conjecture &
an $abc$ inequality with exponent below $4/3$ would exclude all large triples; the exponents derived from the explicit form of the
conjecture are larger, and that form gives only a conditional bound on the middle (Remark~\ref{I:rem:abc}) &
an unconditional route \\
Primitive parts of Lucas sequences (after the primitive divisor theorems of Carmichael \cite{Carmichael1913} and of Bilu, Hanrot and Voutier \cite{BiluHanrotVoutier2001}) &
every divisor $d \ge 5$ of an odd index gives a prime of rank $d$, and every such prime satisfies $p \equiv \pm 1 \pmod{4d}$
(Corollary~\ref{I:cor:teiler}); on \Z{zquadkerne} of the \Z{rgkerne} kernels every building block has at least two distinct prime factors
(Remark~\ref{I:rem:zweiprimitive}) &
whether a building block can be powerful; Problem~\problemvon{I:q:primindex} \\
How many pairs there are (upper bounds at distance two by Chan \cite{Chan2012}, Blomer and Sch{\"o}bel \cite{BlomerSchoebel2013}, and
Tao \cite{Tao2026}) &
an explicit lower bound $\liminf N_1(x)/\log x \ge \Z{csechs}$ at distance one (Proposition~\ref{II:prop:csechs}); at distance two
a bound $N^{\theta}$ for the number of kernels with $m \mid U_1(m)$ and $T_1(m)$ even gives the bound $N^{\max(2/9, \theta)}$ for all
pairs, since the two counts differ by $O(N^{2/9})$ (Proposition~\ref{II:prop:lokalisierung}), and under a squarefreeness assumption the case in which Tao's argument is weakest lies
among those kernels (Remark~\ref{II:rem:tao}) &
the true growth; Problems~\problemvon{II:q:konvergenz}, \problemvon{II:q:wachstum} and~\problemvon{II:q:ezwei} \\
Powerful values of polynomials (De Koninck, Doyon and Luca \cite{DeKoninckDoyonLuca2011}) &
two explicit infinite chains, each giving powerful values of both $x^2 + x + 1$ and $x^2 - x + 1$ (Proposition~\ref{III:prop:phidrei});
one of them contains the prime $x$ reported by Lavi \cite{Lavi2026}; $x^4 + 1 = q^3 a^2$ has no solution for primes $q < 2 \cdot 10^{11}$
(Proposition~\ref{III:prop:phiacht}, by a theorem of Cohn \cite{Cohn1997b} and the published verification of the Ankeny--Artin--Chowla
conjecture \cite{vanderPoortenteRieleWilliams2001,vanderPoortenteRieleWilliams2003}) &
whether $x^4 + 1$ is ever powerful; Problem~\problemvon{III:q:phiacht} \\
Katz's conjecture on Wieferich fractions \cite{Katz2015} &
on \Z{wf3kerne} units and the primes below $\Z{wf3p}$ the Wieferich fractions show no deviation from equidistribution in our measurement, which states its power
(Remark~\ref{III:rem:wieferichbrueche}) &
the conjecture itself \\
\hline
\end{longtable}
\endgroup

\section*{The open problems}\addcontentsline{toc}{section}{The open problems}

The parts end with questions. We collect them here with letters, in the order of the parts, in the wording of the parts. The discussion of
each question, that is, what is known and what an answer would give, stands behind the question in its part. The same list, generated from
the same source, is the file \texttt{OPEN\_PROBLEMS.md} of the data deposit. Erd\H{o}s's own questions, whether a triple exists and whether
$N_1(x)$ is bounded by a power of $\log x$, are not repeated; they are the questions the parts start from.

%
\paragraph{Notation used in the problems.}
\begin{itemize}
\item A positive integer is \emph{powerful} if every prime that divides it divides it at least twice.
\item For a squarefree $m \ge 2$, $\varepsilon_m = T_1 + U_1 \sqrt{m}$ is the fundamental solution of $x^2 - m y^2 = 1$, and $\varepsilon_m^k = T_k + U_k \sqrt{m}$; $T_k(m)$, $U_k(m)$ name the kernel $m$. $m'$ is the product of the primes of $m$ that do not divide $U_1$.
\item The \emph{rank} $\alpha_m(p)$ of a prime $p \nmid m$ is the smallest $k$ with $p \mid T_k(m)$; $M_d$ is the part of $T_d(m)$ supported on the primes of rank exactly $d$ (the \emph{primitive part}).
\item \emph{Wing (ii)} of Part I: the kernels $m \equiv 7 \pmod 8$ with $m \mid U_1(m)$; they are the weak point of the lower bound for a triple.
\item In Part II, the \emph{families} are Walker's Pell families of pairs $n, n + 1$, $F_f$ is the growth factor from one pair of the family $f$ to the next, $w(m) = 1/(2m' \log_{10} \varepsilon_m)$ is the rate of the tower of the kernel $m$, $P_2(N)$ counts the pairs $n - 1, n + 1$ of powerful numbers with $n \le N$, and $E_3(N)$ counts the odd squarefree $m$ with $m \mid U_1(m)$, $T_1(m)$ even and $T_1(m) \le N$.
\item Reinhart's condition (C) is Definition 1.4 of \cite{Reinhart2024}.
\end{itemize}
\par\medskip\noindent\textbf{Problem A} (finiteness of the kernels of wing (ii); Question~\ref{I:q:walsh} in Part~I).\enspace Are there only finitely many squarefree $m \equiv 7 \pmod 8$ with $m \mid U_1(m)$ and $T_1(m)$ even?
\par\medskip\noindent\textbf{Problem B} (the primitive part of $T_d(m)$; Question~\ref{I:q:primindex} in Part~I).\enspace Let $m \equiv 7 \pmod 8$ be squarefree with $T_1(m)$ even, let $d \ge 5$ be odd, and let $\Prim{d}$ be the part of $T_d(m)$ supported on the primes of rank exactly $d$. Can $\Prim{d}$ be powerful?
\par\medskip\noindent\textbf{Problem C} (a middle that is a power of $2$; Question~\ref{I:q:zweierpotenz} in Part~I).\enspace Can $2^a - 1$ and $2^a + 1$ both be powerful for some $a \ge 2$?
\par\medskip\noindent\textbf{Problem D} (Reinhart's condition (C); Question~\ref{I:q:reinhart} in Part~I).\enspace Does Reinhart's condition~(C) \cite[Definition~1.4]{Reinhart2024} hold for some squarefree $d$, that is, is there an index $k$ with $T_k(d)$ powerful and $d \mid U_k(d)$? For $d \mid U_1(d)$ the second condition holds for every $k$.
\par\medskip\noindent\textbf{Problem E} (convergence of the rates; Question~\ref{II:q:konvergenz} in Part~II).\enspace Does $\sum_f 1/\log F_f$, taken over all families, converge?
\par\medskip\noindent\textbf{Problem F} (growth of the tower rate; Question~\ref{II:q:wachstum} in Part~II).\enspace Does $c_M = \sum_{m \le M} w(m)$, the sum over odd squarefree $m$ with $T_1(m)$ even of $w(m) = 1/(2m' \log_{10} \varepsilon_m)$, stay bounded as $M \to \infty$?
\par\medskip\noindent\textbf{Problem G} (the kernels that divide $U_1$; Question~\ref{II:q:ezwei} in Part~II).\enspace Is there a $\theta < 61/180$ with $E_3(N) \ll N^{\theta}$? Since $E_3(N) \le P_2(N)$, Blomer and Sch{\"o}bel give $E_3(N) \ll N^{\gamma}$ for every $\gamma > 61/180$ \cite[Theorem~3]{BlomerSchoebel2013}; the question is whether the squarefree odd $m$ with $m \mid U_1(m)$, $T_1(m)$ even and $T_1(m) \le N$ are rarer than that.
\par\medskip\noindent\textbf{Problem H} (the most frequent gap; Question~\ref{III:q:luecken} in Part~III).\enspace Let $g(x)$ be the most frequent gap between consecutive powerful numbers up to $x$. Is $g(x)$ powerful for all large $x$, and does the share of powerful values among the $N$ most frequent gaps up to $x$ tend to $1$ when $N$ is fixed and $x \to \infty$?
\par\medskip\noindent\textbf{Problem I} (powerful values of $x^4 + 1$; Question~\ref{III:q:phiacht} in Part~III).\enspace Is $x^4 + 1$ powerful for some integer $x \ge 1$?


\medskip
\noindent\textbf{One gap, seen from two sides.} Problems~\problemvon{I:q:walsh} and~\problemvon{II:q:ezwei} concern two nested sets of kernels. Part~I bounds the middle of a
triple from below in two wings (Corollary~\ref{I:cor:fluegel}); wing~(ii) consists of the kernels $m$ with $m \mid U_1(m)$, and it is the
weaker wing. Part~II localizes the count $P_2(N)$ of pairs at distance two in the kernels with $m \mid U_1(m)$ and $T_1(m)$ even, of both
residue classes, of which those with $m \equiv 7 \pmod 8$ form wing~(ii): Proposition~\ref{II:prop:lokalisierung}
says that the count differs from $E_3(N)$ by at most $51 N^{2/9}$, where $E_3(N)$ counts the odd squarefree $m$ with $m \mid U_1(m)$ and $T_1(m)$ even whose first middle is at most $N$. Problem~\problemvon{I:q:walsh} asks whether these
kernels, restricted to $m \equiv 7 \pmod 8$, are finite in number; Problem~\problemvon{II:q:ezwei} asks how fast $E_3(N)$ grows. What is known
is small: Reinhart's list of these kernels is complete up to $1.5 \cdot 10^{12}$ \cite[Remark~5.5]{Reinhart2023} and, if his extended
search missed nothing, up to $5.325 \cdot 10^{13}$ \cite{Reinhart2024} (Remark~\ref{I:rem:reinhartliste}); the number of its odd
kernels with $T_1(m)$ even is \Z{reinhartgeradesicher} below the first bound and \Z{reinhartgerade} below the second (Section~\ref{II:sec:fragen} of Part~II). A bound $E_3(N) \ll N^\theta$ with $\theta < 61/180$ would improve the best exponent
we know for pairs at distance two; a proof that the kernels of wing~(ii) are finite in number would confine that wing to finitely many
kernels, each of which still carries infinitely many admissible indices $k$, and for the known ones Problem~\problemvon{I:q:reinhart} is open.

\medskip
\noindent\textbf{One step behind several problems.} Problems~\problemvon{I:q:primindex} and~\problemvon{I:q:zweierpotenz}, and the witness
statement of Part~I (Theorem~\ref{I:thm:zeugen}), stop at the same step, and so does every route we tried. Whether a prime divides a term
$T_k(m)$ is decided modulo $p$ by the rank of the prime, and this controls all indices at once (Lemma~\ref{I:lem:rang},
Corollary~\ref{I:cor:teiler}). Whether a prime of rank $d$ divides $T_d(m)$ twice is a condition modulo $p^2$ of Wieferich type, and the
structure modulo $p$ does not decide it. Corollary~\ref{I:cor:teiler} turns a counterexample into one such condition at every divisor
$d \ge 5$ of the index, and Remark~\ref{I:rem:aussen} shows that in the outer numbers of a triple no prime can carry exponent one at all.
This describes the reach of the methods used here and of the tools we know; it says nothing about whether a triple exists. As a
reference point from outside our methods, Granville \cite{Granville1986} showed that if there are no three consecutive powerful numbers,
then infinitely many primes $p$ satisfy $p^2 \nmid 2^p - 2$ (Theorem~\ref{I:cor:wieferich}). To our knowledge this last statement is
known only under further assumptions such as the $abc$ conjecture, so a proof of nonexistence would have to include it. We take this as
a sign that the step is hard, not as evidence for either answer.
Problems~\problemvon{I:q:walsh} and~\problemvon{II:q:ezwei} concern the different gap described in the preceding paragraph.

\inwords The questions are printed in one place so that a reader can see them together and check them off. Two of them, Problems~\problemvon{I:q:walsh}
and~\problemvon{II:q:ezwei}, are two questions about nested sets of kernels: how many kernels $m$ divide their own $U_1(m)$ and have an even first term.
Part~I needs to know them, so that its lower bound for a triple has no weak wing; Part~II needs them to be few, so that its count of pairs at
distance two is small. Below $1.5 \cdot 10^{12}$ the number of such kernels is exactly \Z{reinhartgeradesicher}, and whether there are finitely
many is open; Walsh asked it for all kernels dividing $U_1$.
